\documentclass[10pt, a4paper]{article}

% --- Packages ---
\usepackage[utf8]{inputenc}
\usepackage[T1]{fontenc}
\usepackage[margin=0.60in, top=0.72in, bottom=0.72in]{geometry}
\usepackage{xcolor}
\usepackage{titlesec}
\usepackage{booktabs}
\usepackage{tabularx}
\usepackage{longtable}
\usepackage{enumitem}
\usepackage{amsmath, amssymb}
\usepackage{fancyhdr}
\usepackage{microtype}
\usepackage{multicol}
\usepackage{mdframed}
\usepackage[hidelinks]{hyperref}

% --- Monochromatic Color Palette ---
\definecolor{monodark}{HTML}{111111}
\definecolor{monomedium}{HTML}{3A3A3A}
\definecolor{monolight}{HTML}{666666}
\definecolor{monoline}{HTML}{CCCCCC}
\definecolor{monoshade}{HTML}{F0F0F0}
\definecolor{monobox}{HTML}{E8E8E8}

% --- Typography & Spacing ---
\setlist[itemize]{noitemsep, topsep=2pt, leftmargin=14pt}
\setlist[enumerate]{noitemsep, topsep=2pt, leftmargin=14pt}
\emergencystretch=4em

% --- Code Box Style ---
\newmdenv[
  backgroundcolor=monoshade,
  linecolor=monoline,
  linewidth=0.5pt,
  topline=true,
  bottomline=true,
  leftline=true,
  rightline=true,
  innerleftmargin=6pt,
  innerrightmargin=6pt,
  innertopmargin=4pt,
  innerbottommargin=4pt,
  skipabove=4pt,
  skipbelow=4pt,
]{codebox}

% --- Section Headers ---
\titleformat{\section}
  {\color{monodark}\large\bfseries}
  {\thesection.}{0.5em}{}
  [\color{monoline}\titlerule]

\titleformat{\subsection}
  {\color{monomedium}\normalsize\bfseries}
  {\thesubsection}{0.5em}{}

\titleformat{\subsubsection}
  {\color{monolight}\small\bfseries\itshape}
  {}{0em}{}

\titlespacing*{\section}{0pt}{10pt}{4pt}
\titlespacing*{\subsection}{0pt}{7pt}{2pt}
\titlespacing*{\subsubsection}{0pt}{5pt}{1pt}

% --- Header and Footer ---
\pagestyle{fancy}
\fancyhf{}
\renewcommand{\headrulewidth}{0.4pt}
\renewcommand{\footrulewidth}{0.4pt}
\lhead{\small\color{monolight} VIRTUAL AI --- DETAILED TECHNICAL SPECIFICATION}
\rhead{\small\color{monolight} AUGUST 2026}
\lfoot{\small\color{monolight} Om Mukherjee, Aryan Yadav, Aman Vatsa}
\rfoot{\small\color{monolight} Page \thepage}

% =============================================================
\begin{document}
% =============================================================

% ---- TITLE BLOCK ----
\begin{center}
  {\Huge\bfseries\color{monodark} Virtual AI}\\[5pt]
  {\large\itshape\color{monomedium}
    Real-Time AI-Powered Interview Simulator, Evaluator \& NLP Analytics Platform
  }\\[9pt]
  \noindent\rule{\linewidth}{1.2pt}\\[6pt]
  {\small
    \textbf{Authors:} Om Mukherjee \quad Aryan Yadav \quad Aman Vatsa
    \quad\textbullet\quad
    \textbf{Mentors:} Harshala Mahajan \quad Sayali Khandare
  }\\[3pt]
  {\small\color{monolight}
    \url{https://github.com/meekhumor/virtual_interviewer}
    \quad $|$ \quad
    \url{https://usevirtual-ai.vercel.app}
    \quad $|$ \quad
    Backend: \url{https://virtual-ai-iimu.onrender.com}
  }\\[8pt]
  \noindent\rule{\linewidth}{0.5pt}
\end{center}

% ---- ABSTRACT ----
\noindent\textbf{Abstract.}\quad
Virtual AI is a production deployed, full stack AI interview preparation
platform built to simulate realistic technical and behavioural job interviews
through a real-time conversational interface. It combines a React 18 frontend
with a Django REST Framework + Django Channels ASGI backend, a LangGraph
ReAct agent using Google Gemini 2.0 Flash (with automated Groq Llama-3.3-70B
failover), WebRTC video capture, browser Speech-to-Text and Text-to-Speech
synthesis, an embedded Monaco code editor, resume PDF parsing via
\texttt{pdfplumber}, and a comprehensive post-interview NLP analytics engine.
This document provides a full technical specification: system architecture,
API reference, database schema, agent design, authentication mechanisms,
real-time communication, interview lifecycle, frontend state management, and
engineering challenges.

\vspace{6pt}
\noindent\rule{\linewidth}{0.4pt}
\vspace{4pt}

% ---- SWITCH TO TWO-COLUMN BODY ----
\raggedcolumns
\begin{multicols}{2}

% ==========================================
\section{Introduction \& Motivation}
% ==========================================

Technical job interviews test a unique intersection of domain knowledge,
live problem-solving, articulate communication, and composure under
pressure. Candidates frequently cite interview anxiety and lack of realistic
practice as their primary barriers to success.

Traditional mock interview tools fall short in several ways: they rely on
static question banks, lack contextual follow-up, provide no voice
interaction, and fail to deliver personalized feedback grounded in the
candidate's own resume. Virtual AI was designed to close every one of these
gaps.

\subsection{Project Goals}
\begin{enumerate}
  \item Simulate a real human interviewer through an AI agent that listens,
        evaluates, and dynamically follows up.
  \item Support both \textbf{Practice Mode} (constructive, encouraging) and
        \textbf{Real Mode} (neutral, strict panel format).
  \item Generate meaningful, multi-dimensional performance feedback rather
        than a simple pass/fail score.
  \item Support full multimodal input: voice, typed text, and live code.
  \item Record video responses for self-review.
  \item Allow candidate resume upload to personalise question generation.
\end{enumerate}

% ==========================================
\section{Complete Technology Stack}
% ==========================================

\begin{center}
\small\color{monodark}
\begin{tabular}{@{}p{1.8cm} p{4.5cm}@{}}
\toprule
\textbf{Layer} & \textbf{Technology} \\
\midrule
Frontend     & React 18, Vite, Tailwind CSS, React Router v6 \\
UI Libraries & Lucide React, Lottie Web, react-draggable \\
Voice        & Web Speech STT, \texttt{speechSynthesis} TTS \\
Code Editor  & \texttt{@monaco-editor/react} \\
Video        & WebRTC, MediaRecorder API (\texttt{.webm}) \\
WebSocket    & Native \texttt{WebSocket} API (custom hook) \\
Backend      & Python 3.10, Django 5.x, DRF, Channels \\
ASGI Server  & Daphne / Uvicorn \\
AI Framework & LangGraph, LangChain \\
LLMs         & Gemini 2.0 Flash (primary), Groq Llama-3.3-70B \\
Auth         & Supabase Auth (JWT) \\
Database     & PostgreSQL (Supabase-managed) \\
Storage      & Supabase Storage (S3-compatible) \\
PDF Parsing  & \texttt{pdfplumber} \\
Testing      & pytest-django, pytest-asyncio, channels.testing \\
Hosting      & Vercel (frontend), Render (backend) \\
\bottomrule
\end{tabular}
\end{center}

% ==========================================
\section{High-Level System Architecture}
% ==========================================

The system follows a \textit{decoupled client-server} model with two distinct
communication channels operating in parallel:

\subsubsection{Channel 1 --- HTTP REST (stateless)}
All resource management operations (auth, resume upload, interview
creation, analytics retrieval) use authenticated HTTP REST endpoints under
\texttt{/api/}. Requests carry a Supabase JWT in the \texttt{Authorization}
header.

\subsubsection{Channel 2 --- WebSocket (stateful, real-time)}
The live interview conversation uses a persistent WebSocket connection to
\texttt{ws://host/ws/interview/<id>/?token=<jwt>}. The JWT is passed as a
query parameter because the browser \texttt{WebSocket} constructor does not
support custom headers.

\subsubsection{Data Flow Summary}
\begin{codebox}
\footnotesize
\begin{verbatim}
[Browser]
  |-- HTTP POST /api/interviews/start/ -->
  |<- {interview_id, ws_url} ----------
  |
  |-- WebSocket (ws_url?token=...) ---->
  |<- {connected}                 -----
  |<- {agent_response, question}  -----
  |-- {user_message, text, code}  ---->
  |<- {agent_response, feedback}  -----
  ...
  |-- HTTP PATCH /api/interviews/N/status/
  |-- Navigate to /review/N
\end{verbatim}
\end{codebox}

% ==========================================
\section{Authentication System}
% ==========================================

\subsection{Supabase JWT Authentication (REST)}
Class: \texttt{SupabaseJWTAuthentication}
(\texttt{api/auth.py}).

For every protected REST endpoint, the middleware:
\begin{enumerate}
  \item Extracts the Bearer token from the \texttt{Authorization} header.
  \item Calls \texttt{supabase.auth.get\_user(token)} to validate the JWT
        against Supabase's remote auth service.
  \item Looks up the corresponding \texttt{User} record by \texttt{supabase\_id}
        in the local PostgreSQL database.
  \item Returns \texttt{(django\_user, token)} to DRF's permission system.
\end{enumerate}

\subsection{WebSocket JWT Middleware}
Class: \texttt{TokenAuthMiddleware}
(\texttt{api/auth.py}).

ASGI middleware wrapping every WebSocket connection:
\begin{enumerate}
  \item Parses the query string: \texttt{?token=<jwt>}.
  \item Validates via \texttt{sync\_to\_async(supabase.auth.get\_user)}.
  \item Populates \texttt{scope["user"]} before the consumer handles the
        connection.
  \item Sets \texttt{scope["user"] = AnonymousUser()} on any failure,
        causing \texttt{InterviewConsumer} to reject the connection.
\end{enumerate}

\subsection{Guest Session System}
Endpoint: \texttt{POST /api/guest-login/}
\begin{itemize}
  \item Generates a UUID, creates a local \texttt{User} record with email
        \texttt{guest\_<uuid>@virtualai.local} and a matching
        \texttt{supabase\_id}.
  \item Returns token format: \texttt{guest\_token\_<uuid>}.
  \item Auth middleware recognises the \texttt{guest\_token\_} prefix and
        bypasses Supabase validation, looking up the user by
        \texttt{supabase\_id} directly.
  \item \textbf{Cleanup}: On every new guest login, all guest accounts older
        than 12 hours are deleted atomically before creating the new session.
\end{itemize}

% ==========================================
\section{Database Schema}
% ==========================================

\subsection{User Model}
Extends \texttt{AbstractBaseUser} + \texttt{PermissionsMixin}. Login field
is \texttt{email} (not username). Key fields:

\begin{center}
\small
\begin{tabular}{@{}ll@{}}
\toprule
Field & Notes \\
\midrule
\texttt{email} & \texttt{unique=True}, login identifier \\
\texttt{username} & \texttt{unique=True}, display name \\
\texttt{supabase\_id} & UUID linking to Supabase Auth \\
\texttt{profile\_image\_url} & Supabase Storage public URL \\
\texttt{is\_active / is\_staff} & Standard Django flags \\
\texttt{created\_at} & Auto-set on creation \\
\bottomrule
\end{tabular}
\end{center}

\subsection{Resume Model}
One user may have many resumes; the most recently created is always used
for AI context injection.

\begin{center}
\small
\begin{tabular}{@{}ll@{}}
\toprule
Field & Notes \\
\midrule
\texttt{user} & FK $\to$ User, CASCADE \\
\texttt{file\_url} & Supabase Storage public URL \\
\texttt{parsed\_text} & Up to 3,000 chars, extracted by \texttt{pdfplumber} \\
\texttt{created\_at} & Used for ``most recent'' ordering \\
\bottomrule
\end{tabular}
\end{center}

\subsection{Interview Model}
\begin{center}
\small
\begin{tabular}{@{}ll@{}}
\toprule
Field & Notes \\
\midrule
\texttt{user} & FK $\to$ User \\
\texttt{title} & Default ``General Interview'' \\
\texttt{level} & ENTRY / INTERMEDIATE / SENIOR \\
\texttt{mode} & REAL / PRACTICE \\
\texttt{duration\_seconds} & Set at start; updated at end \\
\texttt{status} & PENDING $\to$ IN\_PROGRESS $\to$ COMPLETED \\
\texttt{scheduled\_at} & Timestamp of start call \\
\bottomrule
\end{tabular}
\end{center}

\subsection{Question Model}
Each interview has an ordered sequence of questions, created progressively
throughout the session. \texttt{order} starts at 1 and increments with
each AI follow-up.

\subsection{Response Model}
\begin{center}
\small
\begin{tabular}{@{}ll@{}}
\toprule
Field & Notes \\
\midrule
\texttt{interview / question} & FKs linking the turn \\
\texttt{text} & Candidate's answer text \\
\texttt{ai\_feedback} & 2--3 sentence AI evaluation \\
\texttt{video\_url} & Supabase Storage URL (\texttt{.webm}) \\
\texttt{agent\_metadata} & JSON, stores \texttt{\{"score": 7\}} \\
\bottomrule
\end{tabular}
\end{center}

\subsection{InterviewAnalysis Model}
A cached, one-to-one post-interview report:
\begin{center}
\small
\begin{tabular}{@{}ll@{}}
\toprule
Field & Notes \\
\midrule
\texttt{overall\_score} & Float 0--10 \\
\texttt{communication\_score} & Float 0--10 \\
\texttt{technical\_score} & Float 0--10 \\
\texttt{confidence\_score} & Float 0--10 \\
\texttt{pace\_wpm} & Estimated WPM (float) \\
\texttt{filler\_word\_count} & Integer count \\
\texttt{power\_word\_count} & Integer count \\
\texttt{strengths} & JSON list of strings \\
\texttt{improvements} & JSON list of strings \\
\texttt{summary} & 2--3 sentence prose summary \\
\bottomrule
\end{tabular}
\end{center}

\subsection{Entity Relationship}
\begin{codebox}
\footnotesize
\begin{verbatim}
User --|<-- Resume (many)
User --|<-- Interview (many)
         |-- Question (many, ordered)
         |     |-- Response (one)
         |-- InterviewAnalysis (one)
\end{verbatim}
\end{codebox}

% ==========================================
\section{REST API Endpoint Reference}
% ==========================================

All endpoints are prefixed with \texttt{/api/}.
Protected endpoints require \texttt{Authorization: Bearer <token>}.

\begin{center}
\scriptsize
\begin{tabular}{@{}p{0.4cm} p{3.3cm} p{2.6cm}@{}}
\toprule
\textbf{M} & \textbf{Endpoint} & \textbf{Description} \\
\midrule
POST & \texttt{signup/} & Register; triggers Supabase email confirmation \\
POST & \texttt{login/} & Supabase sign-in, returns JWT \\
POST & \texttt{guest-login/} & Create ephemeral guest session \\
\midrule
GET  & \texttt{profile/} & Return email, username, avatar URL \\
PATCH& \texttt{profile/update/} & Update username \\
POST & \texttt{profile/image/} & Upload avatar to Supabase bucket \\
GET  & \texttt{profile/stats/} & Totals: interviews, time, avg score \\
\midrule
GET  & \texttt{resume/upload/} & Check for existing resume \\
POST & \texttt{resume/upload/} & Upload PDF; parse + cache text \\
\midrule
POST & \texttt{interviews/start/} & Create interview record + first Q \\
GET  & \texttt{interviews/} & List all interviews (ordered by date) \\
GET  & \texttt{interviews/<id>/} & Full detail with Q\&A \\
DEL  & \texttt{interviews/<id>/} & Delete an interview \\
PATCH& \texttt{interviews/<id>/status/} & Update status + duration \\
POST & \texttt{interviews/<id>/analyze/} & Generate + cache analysis \\
GET  & \texttt{interviews/<id>/analysis/} & Retrieve cached analysis \\
POST & \texttt{interviews/<id>/q/<q>/response/} & HTTP fallback submit \\
POST & \texttt{interviews/video/upload/} & Upload \texttt{.webm} clip \\
\midrule
POST & \texttt{contact/} & Log contact inquiry (no auth required) \\
\bottomrule
\end{tabular}
\end{center}

% ==========================================
\section{LLM Utility Layer}
% ==========================================

\subsection{\texttt{\_call\_llm()} (REST Views)}
A thin synchronous helper used by all non-WebSocket LLM calls:
\begin{codebox}
\footnotesize
\begin{verbatim}
def _call_llm(prompt, max_output_tokens=500):
    try:
        return gemini_client.generate_content(
            model="gemini-2.0-flash",
            contents=prompt, ...
        ).text.strip()
    except Exception:
        return groq_client.chat.completions.create(
            model="llama-3.3-70b-versatile",
            messages=[{"role":"user","content":prompt}],
        ).choices[0].message.content.strip()
\end{verbatim}
\end{codebox}
Used by: \texttt{StartInterviewView}, \texttt{SubmitResponseView}
(HTTP fallback), and \texttt{InterviewAnalysisView}.

\subsection{\texttt{\_get\_resume\_text(user)}}
Fetches the \texttt{parsed\_text} of the user's most recent
\texttt{Resume} record, truncated to 3,000 characters.
Called at the start of every LLM prompt to provide candidate context.

% ==========================================
\section{Real-Time WebSocket Agent}
% ==========================================

\subsection{Consumer Lifecycle}
Class: \texttt{InterviewConsumer}
(\texttt{api/consumers.py}).

\textbf{On \texttt{connect()}:}
\begin{enumerate}
  \item Validate interview ownership with \texttt{\_validate\_interview()}.
  \item Accept the WebSocket.
  \item Load interview metadata (level, mode) and resume text.
  \item Instantiate \texttt{ChatGoogleGenerativeAI} (Gemini) and
        \texttt{ChatGroq} (fallback). Both are stored on \texttt{self} for
        per-session persistence.
  \item Build the structured system prompt.
  \item Call \texttt{create\_react\_agent(llm, [save\_interview\_turn],
        prompt=system\_prompt)} to create the LangGraph agent.
  \item Send \texttt{\{"type":"connected"\}} to the client.
  \item If the interview already has a question (e.g., session resumed),
        immediately deliver it; otherwise call
        \texttt{\_generate\_first\_question()}.
\end{enumerate}

\textbf{On \texttt{receive()}:}
\begin{enumerate}
  \item Parse incoming JSON; reject non-\texttt{user\_message} types silently.
  \item Assemble \texttt{agent\_input} with the full message history,
        recent history summary (last 5 turns), level/mode context, and the
        candidate's text.
  \item Call \texttt{\_invoke\_with\_fallback(agent\_input)}.
  \item Update \texttt{self.messages} with the returned message list for
        next-turn context.
  \item If a \texttt{video\_url} was included, update the latest
        \texttt{Response} record asynchronously.
  \item Send the final AI message text as \texttt{\{"type":"agent\_response"\}}.
\end{enumerate}

\textbf{On \texttt{disconnect():}}
Logs the close code for monitoring. No cleanup required as all state was
persisted at each turn by the agent tool.

\subsection{System Prompt Architecture}
The system prompt enforces a strict \textbf{four-step cognitive loop} per
candidate turn:

\textbf{Step 1 --- Evaluate.} Internally assess: technical accuracy,
concept depth, communication clarity, and relevance. Assign score 1--10.
Note one or two specific strengths/weaknesses.

\textbf{Step 2 --- Formulate.} Generate ONE focused follow-up question:
\begin{itemize}
  \item Must build on what the candidate just said.
  \item Must probe a weakness or unexplored concept.
  \item Must match the level: ENTRY = concept understanding; INTERMEDIATE
        = applied problem solving; SENIOR = system design, trade-offs.
\end{itemize}

\textbf{Step 3 --- Act.} Call \texttt{save\_interview\_turn} exactly once
with: \texttt{candidate\_response}, \texttt{feedback}, \texttt{next\_question}
(no preamble), and \texttt{score}.

\textbf{Step 4 --- Speak.} Reply to the candidate in $\le$35 words:
\texttt{[Optional brief acknowledgement]. [Next question]}.

\textbf{Tone:}
\begin{itemize}
  \item REAL mode: professional, neutral, no excessive praise.
  \item PRACTICE mode: warm, honest, constructive, encouraging.
\end{itemize}

\subsection{LangGraph ReAct Tool: \texttt{save\_interview\_turn}}

This is the bridge between the agent and the database. Decorated with
\texttt{@tool} (LangChain), called asynchronously by the agent. It
delegates to \texttt{\_save\_turn\_to\_db()} which wraps all writes in
\texttt{transaction.atomic()}:
\begin{enumerate}
  \item Fetch the current interview's highest-order \texttt{Question}.
  \item Create a \texttt{Response} linked to that question with the
        candidate text, AI feedback, score, and an empty \texttt{video\_url}
        (updated later if video was recorded).
  \item Create the \texttt{next\_question} record at order $+1$.
\end{enumerate}
Both writes are atomic: either both succeed or neither does.

Tool signature:
\begin{codebox}
\footnotesize
\begin{verbatim}
async def save_interview_turn(
    interview_id: str,
    candidate_response: str,
    feedback: str,
    next_question: str,   # question text ONLY
    score: int = None
) -> str
\end{verbatim}
\end{codebox}

\subsection{Zero-Downtime LLM Failover}
Method: \texttt{\_invoke\_with\_fallback(agent\_input)}
\begin{codebox}
\footnotesize
\begin{verbatim}
try:
    return await self.agent.ainvoke(agent_input)
except Exception as e:
    # Rebuild agent with Groq
    fallback_agent = create_react_agent(
        self.llm_fallback,
        [save_interview_turn],
        prompt=self.system_prompt
    )
    result = await fallback_agent.ainvoke(agent_input)
    # Persist for subsequent turns in this session
    self.agent = fallback_agent
    self.llm = self.llm_fallback
    return result
\end{verbatim}
\end{codebox}
This means the user never sees an error if Gemini goes down mid-session;
the swap is transparent and the conversation continues uninterrupted.

% ==========================================
\section{Resume Processing Pipeline}
% ==========================================

\subsubsection{Upload Flow}
\begin{enumerate}
  \item Client sends \texttt{multipart/form-data} to
        \texttt{POST /api/resume/upload/}.
  \item Backend reads file bytes into memory.
  \item Uploads raw bytes to Supabase Storage bucket \texttt{resumes/}
        at path \texttt{resumes/<supabase\_id>/<filename>}.
  \item Retrieves the public URL via \texttt{get\_public\_url()}.
  \item Passes bytes to \texttt{pdfplumber.open(BytesIO(file\_bytes))}.
        Joins extracted text from all pages; truncates to 3,000 characters.
  \item Creates a \texttt{Resume} database record with both the URL and the
        cached plaintext.
\end{enumerate}

\subsubsection{Context Injection}
At interview start and at each WebSocket connection, \texttt{\_get\_resume\_text()}
fetches the most recent \texttt{Resume.parsed\_text} for the user and
injects it directly into the LLM system prompt:

\begin{codebox}
\footnotesize
\begin{verbatim}
"Candidate resume summary: {resume_text[:3000]}"
\end{verbatim}
\end{codebox}

This allows the AI agent to ask targeted questions about the candidate's
actual stack, projects, and experience rather than generic questions.

% ==========================================
\section{Interview Lifecycle (End-to-End)}
% ==========================================

\begin{enumerate}
  \item User selects \textbf{Interview Setting}: level (Internship / Entry
        / Mid / Management) $\to$ mapped to backend enum (ENTRY / ENTRY /
        INTERMEDIATE / SENIOR), duration (20/30/45/60 min), and mode
        (PRACTICE / REAL). Settings stored in \texttt{sessionStorage}.
  \item User uploads or skips resume on \texttt{/resume}.
  \item \texttt{/cam-permission} requests camera and microphone access.
  \item \texttt{/cam-preview1} or \texttt{/cam-preview2} lets the user
        test their setup, including a live mic test via Web Speech STT.
  \item \texttt{/interview-simulator} is loaded. On ``Start'', it calls
        \texttt{POST /api/interviews/start/} which creates the
        \texttt{Interview} record and generates the opening question via
        \texttt{\_call\_llm()}. The response includes \texttt{ws\_url}.
  \item The custom hook \texttt{useInterviewWS(ws\_url, ...)} opens a
        WebSocket with the JWT token. The agent delivers the first question
        over TTS.
  \item For each turn:
    \begin{itemize}
      \item STT transcribes speech in real time.
      \item Silence detection (configurable timeout) or manual submit
            triggers response submission.
      \item In REAL mode, \texttt{MediaRecorder} captures a \texttt{.webm}
            clip, uploaded via \texttt{POST /api/interviews/video/upload/}.
      \item The WebSocket message includes the text and optionally
            \texttt{video\_url} and \texttt{code}.
      \item The agent evaluates, calls the tool, and responds.
    \end{itemize}
  \item When time runs out (countdown reaches 0) or the user clicks
        ``End'', \texttt{endInterview()} is called:
        cancels TTS, stops STT, stops MediaRecorder, closes WebSocket,
        calls \texttt{PATCH .../status/} with \texttt{COMPLETED} and the
        elapsed \texttt{duration\_seconds}, exits fullscreen, and navigates
        to \texttt{/review/<interview\_id>}.
\end{enumerate}

% ==========================================
\section{Frontend Architecture}
% ==========================================

\subsection{Routing (\texttt{App.jsx})}
React Router v6 with a \texttt{<Layout>} wrapper. Protected routes are
gated by \texttt{<ProtectedRoute>} which checks for a valid token in
\texttt{localStorage} and redirects to \texttt{/register} if absent.

Key routes:
\begin{center}
\scriptsize
\begin{tabular}{@{}ll@{}}
\toprule
Path & Component \\
\midrule
\texttt{/} & Animation landing page \\
\texttt{/home} & Full marketing landing page \\
\texttt{/dashboard} & Interview type selection + stats \\
\texttt{/interview-setting} & Level / duration / mode picker \\
\texttt{/resume} & Upload or skip resume \\
\texttt{/cam-permission} & Media permissions gate \\
\texttt{/cam-preview1} & Mic test + camera preview \\
\texttt{/interview-simulator} & Live interview session \\
\texttt{/review-interview} & List of all past interviews \\
\texttt{/review/:id} & Q\&A transcript + video playback \\
\texttt{/analysis} & NLP analytics dashboard \\
\texttt{/profile} & User profile edit \\
\texttt{/domain} & Job-based interview selector \\
\texttt{/courses} & Curated learning resources \\
\bottomrule
\end{tabular}
\end{center}

\subsection{WebSocket Hook: \texttt{useInterviewWS}}
Encapsulates the entire WebSocket lifecycle in a reusable hook:
\begin{itemize}
  \item Opens a \texttt{new WebSocket(url + "?token=" + jwt)} when
        \texttt{wsUrl} becomes non-null.
  \item Dispatches parsed JSON to the correct callback:
        \texttt{onMessage} for \texttt{agent\_response},
        \texttt{onError} for \texttt{error} frames.
  \item Exposes \texttt{sendMessage(obj)}, \texttt{isConnected},
        and \texttt{closeWS()}.
  \item Cleanup closes the socket on unmount or URL change.
\end{itemize}

\subsection{Interview Simulator State Machine}
The component (\texttt{Interview\_Simulator.jsx}, 1028 lines) manages a
complex state machine via \texttt{useState} and \texttt{useEffect}
combinations:

\begin{center}
\scriptsize
\begin{tabular}{@{}ll@{}}
\toprule
State & Purpose \\
\midrule
\texttt{aiSpeaking} & TTS active; blocks STT \\
\texttt{userSpeaking} & STT transcript changing \\
\texttt{waitingForUser} & AI has finished; STT auto-starts \\
\texttt{processing} & Message sent; awaiting agent reply \\
\texttt{isConnected} & WebSocket open \\
\texttt{timeLeft} & Countdown in seconds \\
\texttt{isRunning} & Timer active \\
\texttt{transcriptHistory} & Array of \{sender, text\} turn log \\
\texttt{webcamStream} & \texttt{MediaStream} object \\
\texttt{showCodeEditor} & Monaco drawer toggle \\
\texttt{showChat} & Chat transcript panel toggle \\
\bottomrule
\end{tabular}
\end{center}

\subsubsection{TTS Voice Selection}
\texttt{getPreferredVoice()} iterates the browser's \texttt{speechSynthesis}
voice list in descending preference order:
Google US/UK English $\to$ Samantha/Siri/Daniel $\to$ Microsoft David/Zira
$\to$ any \texttt{en-US}/\texttt{en-GB} voice $\to$ any English voice.

\subsubsection{Silence Detection}
A configurable \texttt{silenceTimeout} (stored in \texttt{sessionStorage})
triggers auto-submit after the transcript stops changing. Setting it to 0
disables auto-submit (manual mode). The timer resets on each transcript
update.

\subsubsection{WebRTC Video Recording (REAL Mode)}
\begin{enumerate}
  \item \texttt{getUserMedia(\{video: true, audio: true\})} when
        \texttt{waitingForUser} becomes true.
  \item \texttt{MediaRecorder} is started with
        \texttt{video/webm;codecs=vp9,opus}, falling back to the browser
        default.
  \item Chunks are pushed to \texttt{recordedChunksRef} every 1 second.
  \item On submit, \texttt{stopMediaRecorder()} assembles a \texttt{Blob},
        which is uploaded to \texttt{POST /api/interviews/video/upload/}.
  \item The returned Supabase URL is sent in the WebSocket message.
\end{enumerate}

\subsection{Visual Feedback Components}
\subsubsection{\texttt{WaveBar}}
Five animated vertical bars that oscillate between 4px and 22px height
using a keyframe animation (\texttt{waveBar}). Each bar has a staggered
\texttt{animation-delay} of $i \times 0.12$s. Inactive bars are rendered
as flat 6px stubs in grey.

\subsubsection{\texttt{SpeakOrb}}
A circular orb with a border that changes colour depending on the active
state (\texttt{active ? colour : \#3f3f46}), housing the animated
\texttt{WaveBar} and the speaker icon. Used for both the AI and User
speaking indicators.

% ==========================================
\section{Post-Interview NLP Analytics}
% ==========================================

\subsection{Trigger Logic (Frontend)}
\texttt{Analysis.jsx} first attempts a \texttt{GET
/api/interviews/<id>/analysis/} to retrieve any cached report. If a 404
is returned (no report yet), it automatically issues a
\texttt{POST /api/interviews/<id>/analyze/} to generate one.

\subsection{LLM Analysis Prompt}
The backend assembles the full interview transcript:
\begin{codebox}
\footnotesize
\begin{verbatim}
Q1: <question>
A1: <candidate response>
Q2: <question>
A2: <candidate response>
...
\end{verbatim}
\end{codebox}
The prompt instructs the LLM to return a strict JSON object with all
metric fields. The response is stripped of markdown fences via regex before
\texttt{json.loads()}.

\subsection{Scoring Formula}
\begin{equation}
\text{overall\_score} = f\!\left(\text{tech}, \text{comm}, \text{conf},
    \text{pace}, \text{fillers}, \text{power words}\right)
\end{equation}
The exact weighting is determined by the LLM's qualitative evaluation.
All numeric scores are on a 0--10 float scale.

\subsection{Pace Calculation}
\begin{equation}
\text{WPM} = \frac{\sum_i \left|\text{tokens in } A_i\right|}{T_{\text{elapsed}} / 60}
\end{equation}

\subsection{Improvement Tips Engine}
\texttt{Analysis.jsx} embeds a hard-coded tip library keyed by metric name
(\texttt{"Pace"}, \texttt{"Filler Words"}, \texttt{"Power Words"},
\texttt{"Communication"}, \texttt{"Technical"}, \texttt{"Confidence"}).
Clicking a metric card opens a modal with 5 specific, actionable tips.

\subsection{Inline Transcript Analyser (\texttt{Check.jsx})}
A standalone developer/debug utility available at \texttt{/check}:
counts filler words, power words, and negative sentiment words in a pasted
transcript, and computes WPM given a manual duration input. This was
used during development to validate the analytics logic.

% ==========================================
\section{Serializers \& Data Contracts}
% ==========================================

\subsubsection{\texttt{QuestionSerializer}}
Uses a \texttt{SerializerMethodField} to attach the related
\texttt{Response} to each question. The \texttt{interview} context
is passed at instantiation time; if absent, the response field returns
\texttt{null}.

\subsubsection{\texttt{InterviewDetailSerializer}}
Overrides \texttt{get\_questions} to force ordering by
\texttt{question.order}, guaranteeing the Q\&A sequence is always
chronological regardless of insertion order.

\subsubsection{\texttt{InterviewAnalysisSerializer}}
Exposes all 12 metric fields of the \texttt{InterviewAnalysis} model
as a flat JSON payload. The \texttt{strengths} and \texttt{improvements}
fields are \texttt{JSONField} lists serialised natively.

% ==========================================
\section{Key Engineering Challenges}
% ==========================================

\subsection{WebSocket JWT Without Custom Headers}
The browser's \texttt{WebSocket} constructor does not support custom HTTP
headers (unlike \texttt{fetch}). The JWT cannot be sent in
\texttt{Authorization}. \textbf{Solution}: pass the token as a URL query
parameter \texttt{?token=<jwt>}. The custom ASGI
\texttt{TokenAuthMiddleware} intercepts the scope before the consumer
is invoked, validates the token, and attaches the \texttt{User} to
\texttt{scope["user"]}.

\subsection{LLM JSON Reliability}
LLMs occasionally prefix or wrap their JSON output in markdown fences
or prose. \textbf{Solution}: post-process every LLM response with:
\begin{codebox}
\footnotesize
\begin{verbatim}
raw = re.sub(r"^```(?:json)?\n?", "", raw)
raw = re.sub(r"```$", "", raw).strip()
data = json.loads(raw)
\end{verbatim}
\end{codebox}

\subsection{Race Condition: Video Upload vs.\ WebSocket Message}
In REAL mode, video must be uploaded \textit{before} the WebSocket
message is sent (since the message needs the URL). But the upload is
async and may take time. \textbf{Solution}: \texttt{submitUserResponse()}
is \texttt{async}; it \texttt{await}s both \texttt{stopMediaRecorder()} and
\texttt{uploadVideoBlob()} before calling \texttt{handleSendMessage()}.
If upload fails, it gracefully sends the message with an empty
\texttt{video\_url}.

\subsection{Resume Token Budget}
A full CV can easily exceed the context window of smaller LLMs.
\textbf{Solution}: \texttt{pdfplumber} extracts only the text layer of
the PDF at upload time and the backend hard-caps at 3,000 characters.

\subsection{STT Stopping During TTS}
Without coordination, the microphone picks up the AI's synthesised voice
as user speech. \textbf{Solution}: \texttt{speakText()} calls
\texttt{SpeechRecognition.stopListening()} in \texttt{utt.onstart} and
only calls \texttt{startListening()} again in \texttt{utt.onend}, after
a 300ms cooldown. The \texttt{aiSpeaking} state flag additionally gates
the auto-start \texttt{useEffect}.

\subsection{Concurrent DB Sessions During Tests}
The remote Supabase PostgreSQL instance had other connections open when
pytest-django tried to tear down the test database
(\texttt{test\_postgres}). This generated a
\texttt{PytestWarning: OperationalError: database is being accessed by
other users}. This is non-blocking (tests still pass) and is a known
artefact of using a shared remote DB for testing.

% ==========================================
\section{Testing Strategy}
% ==========================================

\subsection{Test Suite Location \& Configuration}
File: \texttt{api/tests.py} (291 lines).
Runner: \texttt{venv/bin/pytest} with \texttt{backend.settings} as the
Django settings module.

\subsection{Test Cases}

\begin{center}
\scriptsize
\begin{tabular}{@{}p{3.8cm} p{2.5cm}@{}}
\toprule
\textbf{Test} & \textbf{Validates} \\
\midrule
\texttt{unauthenticated\_rejected} & AnonymousUser cannot connect \\
\texttt{authenticated\_accepted} & Valid user connects, gets Q \\
\texttt{wrong\_user\_rejected} & User cannot access others' interview \\
\texttt{stays\_alive\_after\_q} & WS not closed after first question \\
\texttt{not\_close\_between\_turns} & WS open between AI and user turns \\
\texttt{multiple\_turns} & Two full turn exchanges succeed \\
\texttt{malformed\_json\_error} & Non-JSON input returns error frame \\
\texttt{unknown\_type\_ignored} & Unknown message type silently dropped \\
\texttt{empty\_message\_ignored} & Whitespace text produces no response \\
\texttt{clean\_disconnect} & Client close does not crash server \\
\texttt{disconnect\_is\_logged} & Consumer logs disconnection at INFO \\
\bottomrule
\end{tabular}
\end{center}

\textbf{Mocking strategy}: \texttt{ChatGoogleGenerativeAI},
\texttt{ChatGroq}, and \texttt{create\_react\_agent} are all stubbed with
\texttt{unittest.mock.MagicMock} + \texttt{AsyncMock} to avoid real API
calls. Authentication middleware is bypassed by patching
\texttt{TokenAuthMiddleware.\_\_call\_\_} to inject a real database user.

All 11 tests pass with \texttt{comm.connect(timeout=5)} to allow for
remote PostgreSQL roundtrip latency.

% ==========================================
\section{Deployment Architecture}
% ==========================================

\begin{center}
\small
\begin{tabular}{@{}ll@{}}
\toprule
\textbf{Service} & \textbf{Platform} \\
\midrule
Frontend (React/Vite build) & Vercel \\
Backend (Django/Daphne) & Render \\
PostgreSQL database & Supabase managed \\
File storage & Supabase Storage \\
Auth provider & Supabase Auth \\
\bottomrule
\end{tabular}
\end{center}

The Render backend uses Daphne to serve ASGI, supporting both HTTP and
WebSocket traffic from a single process. \texttt{RENDER\_EXTERNAL\_HOSTNAME}
is auto-appended to \texttt{ALLOWED\_HOSTS} at startup. Static files are
served by WhiteNoise without a separate CDN. WebSocket URLs auto-detect
\texttt{https://} vs.\ \texttt{http://} and emit \texttt{wss://} vs.\
\texttt{ws://} accordingly.

\subsection{Environment Variables Required}
\begin{codebox}
\footnotesize
\begin{verbatim}
DJANGO_SECRET_KEY=<key>
DJANGO_DEBUG=False
SUPABASE_URL=https://<project>.supabase.co
SUPABASE_SECRET_KEY=<key>
GEMINI_API_KEY=<key>
GROQ_API_KEY=<key>
DB_USER=<supabase_db_user>
DB_PASSWORD=<supabase_db_pass>
DB_HOST=<supabase_pooler_host>
DB_NAME=postgres
FRONTEND_URL=https://usevirtual-ai.vercel.app
\end{verbatim}
\end{codebox}

% ==========================================
\section{Future Work}
% ==========================================

\begin{enumerate}
  \item \textbf{Live Facial Emotion Analysis}: Integrate OpenCV or a cloud
        vision API to detect confidence, engagement, and stress from the
        real-time video feed.
  \item \textbf{AI Avatar Interviewer}: Render a photorealistic animated
        avatar using Heygen or similar, lip-synced to the TTS output.
  \item \textbf{Job Portal Integration}: Connect directly to job listings;
        auto-generate tailored question sets from job description text.
  \item \textbf{Custom Question Banks}: Let users or organisations upload
        their own question sets and assign interviews to candidates.
  \item \textbf{Streaming Token Output}: Switch from full-response to
        streaming TTS (word-by-word) to reduce perceived latency.
  \item \textbf{Voice Cloning}: Allow users to practice with a mock
        interviewer voice that matches a real company's style.
\end{enumerate}

% ==========================================
\section{Conclusion}
% ==========================================

Virtual AI is a comprehensive, production-grade demonstration of
full-stack AI engineering. It integrates real-time conversational AI,
multimodal input, persistent ASGI WebSocket state, Supabase cloud
services, automated LLM failover, and a robust NLP analytics pipeline
into a cohesive product that delivers genuine value to job seekers.

The architectural decisions made throughout the project --- decoupled
WebSocket and REST channels, atomic agent tool execution, resume-aware
dynamic prompting, and zero-downtime LLM failover --- each solved concrete
engineering problems and can be applied as reusable patterns in any
real-time AI application.

\end{multicols}
\end{document}
